% ============================================================================
%  ESTADÍSTICA INFERENCIAL — Sesión 14
%  Semana 14: 9 al 14 de noviembre de 2026
%  3.3 Planteamiento y pasos para solucionar una prueba de hipótesis
%  3.4 Uso de software estadístico
% ============================================================================
\documentclass[aspectratio=169,10pt]{beamer}

\usepackage[utf8]{inputenc}
\IfFileExists{spanish.ldf}{\usepackage[spanish,es-noshorthands]{babel}}{}
\usepackage{amsmath,amssymb}
\usepackage{booktabs}
\usepackage{tikz}
\usepackage{listings}

\usetheme{Uniminuto}

\uniasignatura{Estadística inferencial}
\uniprofesor{Andrés Felipe Puerta Velez}
\unicorreo{andres.puerta-v@uniminuto.edu.co}
\uninrc{CBASBL061}
\unisemana{Sesión 14 — Semana 14 · 9 al 14 de noviembre de 2026}

\newcommand{\listatiny}{\setbeamerfont{itemize/enumerate body}{size=\tiny}\tiny}


\newcommand{\fx}[1]{{\ttfamily\color{uniAzul}#1}}

% ------------------------------------------------------------- CÓDIGO R ----
\lstdefinestyle{Rcode}{
  language=R,
  basicstyle=\ttfamily\scriptsize,
  keywordstyle=\color{uniAzul}\bfseries,
  commentstyle=\color{uniGris}\itshape,
  stringstyle=\color{uniTitulo},
  showstringspaces=false, numbers=none, breaklines=true,
  columns=fullflexible,
  backgroundcolor=\color{uniCremaClaro},
  frame=leftline, framerule=1.2pt, rulecolor=\color{uniNaranja},
  xleftmargin=0.35cm, framexleftmargin=0.35cm,
  framexrightmargin=0.15cm, framextopmargin=2pt, framexbottommargin=2pt,
  aboveskip=0.15cm, belowskip=0.15cm
}
\lstdefinestyle{Rmini}{
  language=R,
  basicstyle=\ttfamily\tiny,
  keywordstyle=\color{uniAzul}\bfseries,
  commentstyle=\color{uniGris}\itshape,
  stringstyle=\color{uniTitulo},
  showstringspaces=false, numbers=none, breaklines=true,
  columns=fullflexible,
  backgroundcolor=\color{uniCremaClaro},
  frame=leftline, framerule=1.2pt, rulecolor=\color{uniNaranja},
  xleftmargin=0.35cm, framexleftmargin=0.35cm,
  framexrightmargin=0.15cm, framextopmargin=2pt, framexbottommargin=2pt,
  aboveskip=0.12cm, belowskip=0.12cm
}
\lstdefinestyle{Rout}{
  basicstyle=\ttfamily\tiny\color{black},
  showstringspaces=false, numbers=none, breaklines=true,
  columns=fullflexible,
  backgroundcolor=\color{uniGris!10},
  frame=leftline, framerule=1.2pt, rulecolor=\color{uniGris},
  xleftmargin=0.35cm, framexleftmargin=0.35cm,
  framexrightmargin=0.15cm, framextopmargin=2pt, framexbottommargin=2pt,
  aboveskip=0.12cm, belowskip=0.12cm
}

\begin{document}

\portada

\begin{frame}{Orden del día}
  \begin{enumerate}
    \item Llamado a lista y revisión de los ejercicios propuestos
    \item Repaso: dónde quedamos
    \item \azul{Tema 3.3: los cinco pasos de una prueba de hipótesis}
    \item Pruebas para la media y para la proporción
    \item \azul{Tema 3.4: uso de software estadístico}
    \item Ejercicios en clase y cierre
  \end{enumerate}
\end{frame}

\begin{frame}{Repaso: dónde quedamos}
  En la sesión 13 vimos \azul{qué} es una prueba de hipótesis: $H_0$ y $H_1$,
  errores tipo I y II, $\alpha$ y el valor $p$, con la regla
  \[ p\le\alpha \;\Rightarrow\; \text{se rechaza } H_0 \]

  \vspace{0.2cm}
  \begin{alertblock}{Hoy}
    \footnotesize
    Vemos \destacar{cómo} se ejecuta, paso a paso y con números, y cómo se lee
    la salida de un software cuando el cálculo lo hace la máquina.
  \end{alertblock}
\end{frame}

\seccion{3.3 Planteamiento y pasos}

\begin{frame}{Los cinco pasos}
  \small
  \begin{enumerate}
    \item \azul{Plantear} $H_0$ y $H_1$, e identificar si es unilateral o
          bilateral.
    \item \azul{Fijar} el nivel de significancia $\alpha$ y determinar el
          estadístico de prueba ($z$ o $t$).
    \item \azul{Definir} la región de rechazo (valor crítico) o decidir que se
          usará el valor $p$.
    \item \azul{Calcular} el estadístico con los datos de la muestra.
    \item \azul{Decidir} y \azul{redactar} la conclusión en el contexto del
          problema.
  \end{enumerate}

  \vspace{0.25cm}
  \begin{alertblock}{Dos caminos equivalentes}
    \footnotesize
    \destacar{Valor crítico:} se rechaza si el estadístico cae en la región de
    rechazo. \quad
    \destacar{Valor $p$:} se rechaza si $p\le\alpha$. Siempre llevan a la misma
    decisión.
  \end{alertblock}
\end{frame}

\begin{frame}{Estadísticos de prueba}
  \footnotesize
  \begin{center}
  \begin{tabular}{@{}p{3.4cm}p{3.6cm}p{4.2cm}@{}}
    \toprule
    \textbf{Parámetro} & \textbf{Condición} & \textbf{Estadístico} \\
    \midrule
    Media, $\sigma$ conocida & $n\ge 30$ o población normal &
      $z=\dfrac{\bar{x}-\mu_0}{\sigma/\sqrt{n}}$ \\[0.8em]
    Media, $\sigma$ desconocida & población aprox. normal &
      $t=\dfrac{\bar{x}-\mu_0}{s/\sqrt{n}}$, \; $gl=n-1$ \\[0.8em]
    Proporción & $np_0\ge 5$ y $n(1-p_0)\ge 5$ &
      $z=\dfrac{\hat{p}-p_0}{\sqrt{p_0(1-p_0)/n}}$ \\
    \bottomrule
  \end{tabular}
  \end{center}

  \vspace{0.2cm}
  \begin{alertblock}{Detalle que se olvida}
    \footnotesize
    En la prueba de proporción el error estándar se calcula con
    \destacar{$p_0$}, el valor de $H_0$, no con $\hat p$. En el intervalo de
    confianza sí se usaba $\hat p$, porque allí no hay ningún $p_0$ supuesto.
  \end{alertblock}
\end{frame}

\begin{frame}{Ejercicio 1 — en clase}
  \begin{alertblock}{Individual (10 minutos)}
    El tiempo promedio histórico de atención es de \destacar{15 minutos} con
    $\sigma=2{,}4$. Tras un cambio de proceso, una muestra de
    \destacar{$n=36$} servicios da $\bar{x}=14{,}2$ minutos. Se quiere probar
    si el proceso \destacar{redujo} el tiempo, con $\alpha=0{,}05$.

    \vspace{0.25cm}
    \begin{enumerate}
      \item Plantee $H_0$ y $H_1$ e indique el tipo de prueba.
      \item Calcule el estadístico de prueba.
      \item Compare con el valor crítico y decida.
      \item Redacte la conclusión.
    \end{enumerate}
  \end{alertblock}
\end{frame}

\begin{frame}{Ejercicio 1 — solución}
  \small
  \begin{block}{Pasos 1 a 4}
    \footnotesize
    \begin{enumerate}
      \item $H_0:\mu\ge 15$ \quad vs. \quad $H_1:\mu<15$. Unilateral
            izquierda.
      \item $\sigma$ es conocida y $n=36\ge 30$: se usa $z$.
      \item Valor crítico: $z_{0{,}05}=\azul{-1{,}645}$. Se rechaza si
            $z<-1{,}645$.
      \item \[ z=\frac{14{,}2-15}{2{,}4/\sqrt{36}}
                =\frac{-0{,}8}{0{,}40}=\azul{-2{,}00}
             \qquad p=\azul{0{,}0228} \]
    \end{enumerate}
  \end{block}

  \vspace{0.1cm}
  \begin{exampleblock}{Paso 5}
    \footnotesize
    Como $-2{,}00<-1{,}645$ (y $0{,}0228\le 0{,}05$), \destacar{se rechaza
    $H_0$}. Con un nivel de significancia del 5\%, hay evidencia suficiente
    para concluir que el nuevo proceso redujo el tiempo promedio de atención.
  \end{exampleblock}
\end{frame}

\begin{frame}{Ejercicio 2 — en clase}
  \begin{alertblock}{En parejas (10 minutos)}
    Un proveedor asegura que el contenido promedio de sus empaques es de
    \destacar{500 g}. Una muestra de \destacar{$n=25$} empaques da
    $\bar{x}=512$ g y $s=28$ g. Se quiere verificar si el contenido
    \destacar{difiere} de 500 g, con $\alpha=0{,}05$.

    \vspace{0.25cm}
    \begin{enumerate}
      \item Plantee las hipótesis e indique el tipo.
      \item ¿Se usa $z$ o $t$? ¿Con cuántos grados de libertad?
      \item Calcule el estadístico y compárelo con el valor crítico
            $t_{0{,}025;24}=2{,}064$.
      \item Decida y redacte la conclusión.
    \end{enumerate}
  \end{alertblock}
\end{frame}

\begin{frame}{Ejercicio 2 — solución}
  \small
  \begin{block}{Desarrollo}
    \footnotesize
    $H_0:\mu=500$ \; vs. \; $H_1:\mu\neq 500$. Bilateral. Como el dato es $s$,
    se usa $t$ con $gl=24$.
    \[ t=\frac{512-500}{28/\sqrt{25}}=\frac{12}{5{,}6}=\azul{2{,}1429}
       \qquad p=\azul{0{,}0425} \]
  \end{block}

  \vspace{0.1cm}
  \begin{exampleblock}{Decisión}
    \footnotesize
    $|2{,}1429|>2{,}064$, y además $0{,}0425\le 0{,}05$: \destacar{se rechaza
    $H_0$}. Al 5\% de significancia hay evidencia de que el contenido
    promedio difiere de los 500 g declarados.
  \end{exampleblock}

  \vspace{0.1cm}
  \begin{alertblock}{Observación}
    \footnotesize
    Con $\alpha=0{,}01$ la decisión se invertiría: $0{,}0425>0{,}01$. El
    resultado está \azul{en el límite}, algo que conviene reportar en vez de
    ocultar.
  \end{alertblock}
\end{frame}

\begin{frame}{Ejercicio 3 — en clase}
  \begin{alertblock}{Individual (9 minutos)}
    La meta de satisfacción es del \destacar{60\%}. En una encuesta a
    \destacar{$n=400$} clientes, 248 se declaran satisfechos
    ($\hat{p}=0{,}62$). Se quiere probar si la satisfacción
    \destacar{supera} la meta, con $\alpha=0{,}05$.

    \vspace{0.25cm}
    \begin{enumerate}
      \item Plantee las hipótesis.
      \item Calcule el estadístico usando $p_0=0{,}60$.
      \item Decida con el valor crítico $z_{0{,}05}=1{,}645$.
      \item ¿Coincide con la conclusión del intervalo de confianza de la
            sesión 11?
    \end{enumerate}
  \end{alertblock}
\end{frame}

\begin{frame}{Ejercicio 3 — solución}
  \small
  \begin{block}{Desarrollo}
    \footnotesize
    $H_0:p\le 0{,}60$ \; vs. \; $H_1:p>0{,}60$. Unilateral derecha.
    \[ z=\frac{0{,}62-0{,}60}{\sqrt{\dfrac{0{,}60(0{,}40)}{400}}}
        =\frac{0{,}02}{0{,}0245}=\azul{0{,}8165}
       \qquad p=\azul{0{,}2071} \]
  \end{block}

  \vspace{0.1cm}
  \begin{alertblock}{Decisión y punto 4}
    \footnotesize
    $0{,}8165<1{,}645$ y $0{,}2071>0{,}05$: \destacar{no se rechaza $H_0$}. No
    hay evidencia suficiente para afirmar que la satisfacción supere el 60\%.

    \vspace{0.15cm}
    En la sesión 11 el intervalo del 95\% $(0{,}5724;\ 0{,}6676)$
    \azul{contenía} el valor 0{,}60. Ambos métodos coinciden: el 0{,}60 sigue
    siendo un valor plausible.
  \end{alertblock}
\end{frame}

\seccion{3.4 Uso de software estadístico}

\begin{frame}[fragile]{Las cinco pruebas de la unidad, en R}
  \scriptsize
\begin{lstlisting}[style=Rmini]
# 1. Prueba t para una media (sigma desconocida)
t.test(contenido, mu = 500)                       # bilateral
t.test(contenido, mu = 500, alternative = "greater")   # unilateral derecha

# 2. Prueba para una proporcion
prop.test(248, 400, p = 0.60, alternative = "greater")
binom.test(248, 400, p = 0.60, alternative = "greater")  # exacta

# 3. Comparacion de dos medias independientes
t.test(ventas ~ sucursal)

# 4. Comparacion de dos proporciones
prop.test(c(120, 96), c(200, 200))

# 5. Valores criticos y valores p, a mano
qnorm(0.95);  qt(0.975, df = 24)
2 * (1 - pt(2.1429, df = 24))      # valor p bilateral
\end{lstlisting}

  \vspace{-0.12cm}
  {\tiny\color{uniGris} \fx{alternative} acepta \fx{"two.sided"} (predeterminado),
  \fx{"greater"} y \fx{"less"}. \destacar{Elegirlo mal es el error más costoso}:
  cambia el valor $p$ y puede invertir la decisión.}
\end{frame}

\begin{frame}[fragile]{Cómo se lee una salida de R}
  \scriptsize
  Los programas hacen los pasos 3 y 4. Los pasos 1, 2 y 5 siguen siendo suyos.
\begin{lstlisting}[style=Rout]
	One Sample t-test

data:  contenido
t = 2.1429, df = 24, p-value = 0.0425
alternative hypothesis: true mean is not equal to 500
95 percent confidence interval:
 500.4422 523.5578
sample estimates:
mean of x
      512
\end{lstlisting}

  \vspace{-0.05cm}
  \begin{alertblock}{Qué mirar, y en qué orden}
    \tiny
    (1) \fx{alternative}: verificar que R hizo la prueba que usted planteó. \quad
    (2) \fx{df = 24}: confirma que $n = 25$. \quad
    (3) \fx{p-value = 0.0425} $< 0{,}05$: \destacar{se rechaza $H_0$}. \quad
    (4) El intervalo $(500{,}44;\ 523{,}58)$ \destacar{no contiene 500}, lo que
    confirma el rechazo: prueba e intervalo siempre coinciden.
  \end{alertblock}
\end{frame}

\begin{frame}[fragile]{De la salida a la conclusión escrita}
  \scriptsize
\begin{lstlisting}[style=Rmini]
prueba <- t.test(contenido, mu = 500)
prueba$statistic     # el estadistico t
prueba$parameter     # los grados de libertad
prueba$p.value       # el valor p
prueba$conf.int      # el intervalo de confianza
prueba$estimate      # la media muestral
\end{lstlisting}

  \vspace{0.05cm}
  \begin{exampleblock}{El párrafo que se entrega}
    \tiny
    ``Con una muestra de 25 envases se obtuvo un contenido medio de
    \azul{512 ml} ($s = 28$). La prueba $t$ para una media rechaza
    $H_0: \mu = 500$ ($t = 2{,}14$; $gl = 24$; $p = 0{,}0425$) al nivel del
    5\%. El intervalo de confianza del 95\% para el contenido medio es
    \azul{(500,44; 523,56) ml}, de modo que hay evidencia de que el llenado
    supera los 500 ml declarados.''
  \end{exampleblock}

  \vspace{0.05cm}
  \begin{alertblock}{Los tres elementos que nunca faltan}
    \tiny
    El \destacar{estadístico con sus grados de libertad}, el
    \destacar{valor $p$} y el \destacar{tamaño del efecto} —aquí, el intervalo
    y la media—. Un valor $p$ solo no dice si la diferencia importa en la
    práctica.
  \end{alertblock}
\end{frame}

\begin{frame}{Lo que el software no hace}
  \scriptsize
  \begin{center}
  \begin{tabular}{@{}p{4.6cm}p{7.2cm}@{}}
    \toprule
    \textbf{Decisión} & \textbf{Por qué sigue siendo suya} \\
    \midrule
    Plantear $H_0$ y $H_1$ &
      R prueba lo que usted le escriba, aunque no tenga sentido \\
    \addlinespace
    Elegir la prueba adecuada &
      \fx{t.test()} corre igual con datos que no cumplen sus supuestos \\
    \addlinespace
    Fijar $\alpha$ antes de ver los datos &
      Ajustar $\alpha$ después de ver el $p$ invalida la prueba \\
    \addlinespace
    Juzgar si la muestra es probabilística &
      Con datos mal recogidos, R devuelve un $p$ igual de convincente \\
    \addlinespace
    Decidir si la diferencia importa &
      Con $n$ grande, diferencias irrelevantes salen significativas \\
    \bottomrule
  \end{tabular}
  \end{center}

  \vspace{0.05cm}
  \begin{alertblock}{La frase que resume la sesión}
    \tiny
    \destacar{R calcula el valor $p$; usted decide si la pregunta valía la
    pena.} Por eso el examen se resuelve a mano: para que la herramienta sea
    un acelerador y no una excusa.
  \end{alertblock}
\end{frame}

\begin{frame}{Cierre de la sesión}
  \begin{block}{Lo que debe quedar claro hoy}
    \begin{itemize}
      \item Los cinco pasos: plantear, fijar $\alpha$, definir región, calcular,
            decidir y redactar.
      \item $z$ si el dato es $\sigma$; $t$ con $gl=n-1$ si el dato es $s$.
      \item En proporciones el error estándar se calcula con $p_0$.
      \item Valor crítico y valor $p$ llevan siempre a la misma decisión.
      \item El software calcula, pero no piensa por uno.
    \end{itemize}
  \end{block}

  \vspace{0.2cm}
  \begin{alertblock}{Trabajo independiente para la próxima sesión}
    \footnotesize
    Resolver los \destacar{20 ejercicios de la lámina siguiente}. La sesión 15
    cierra la unidad con \azul{correlación y regresión lineal}, y hay repaso
    para el Examen 3.
  \end{alertblock}
\end{frame}

\begin{frame}{Ejercicios propuestos}
  \vspace{-0.1cm}
  \begin{beamercolorbox}[wd=\textwidth,sep=0.10cm]{block body}
    \centering\scriptsize
    $z$ crítico: $\pm1{,}96$ (bilat. 5\%), $1{,}645$ (unilat. 5\%). \quad
    $t_{0{,}025;24}=2{,}064$, \; $t_{0{,}025;19}=2{,}093$.
  \end{beamercolorbox}
  \listatiny
  \vspace{0.15cm}
  \begin{columns}[T,onlytextwidth]
    \begin{column}{0.47\textwidth}
      \begin{enumerate}\setlength{\itemsep}{0.16cm}
        \item Primer paso de una prueba de hipótesis
        \item ¿Quién decide $\alpha$ y cuándo?
        \item Estadístico si $\sigma$ es conocida
        \item Estadístico si el dato es $s$
        \item Estadístico para una proporción
        \item ¿Con $p_0$ o con $\hat p$ el error estándar?
        \item $\bar x=14{,}2$, $\mu_0=15$, $\sigma=2{,}4$, $n=36$: valor de $z$
        \item Valor $p$ del punto 7 (cola izquierda)
        \item Decisión del punto 7 con $\alpha=0{,}05$
        \item $\bar x=512$, $\mu_0=500$, $s=28$, $n=25$: valor de $t$
      \end{enumerate}
    \end{column}
    \begin{column}{0.50\textwidth}
      \begin{enumerate}\setlength{\itemsep}{0.16cm}
        \setcounter{enumi}{10}
        \item Grados de libertad del punto 10
        \item Decisión del punto 10 con $\alpha=0{,}05$
        \item Decisión del punto 10 con $\alpha=0{,}01$
        \item $\hat p=0{,}62$, $p_0=0{,}60$, $n=400$: valor de $z$
        \item Decisión del punto 14 con $\alpha=0{,}05$
        \item ¿Coincide con el IC de la sesión 11?
        \item ¿Qué se mira primero en una salida de software?
        \item ¿El software verifica si la prueba es la adecuada?
        \item Si $p=0{,}07$ y $\alpha=0{,}10$: decisión
        \item ¿Rechazar $H_0$ prueba que $H_1$ es cierta?
      \end{enumerate}
    \end{column}
  \end{columns}
\end{frame}

\begin{frame}{Respuestas}{Autoevaluación de los ejercicios propuestos}
  \listatiny
  \vspace{0.15cm}
  \begin{columns}[T,onlytextwidth]
    \begin{column}{0.47\textwidth}
      \begin{enumerate}\setlength{\itemsep}{0.16cm}
        \item Plantear $H_0$ y $H_1$
        \item El investigador, antes de calcular
        \item $z=(\bar x-\mu_0)/(\sigma/\sqrt n)$
        \item $t=(\bar x-\mu_0)/(s/\sqrt n)$, $gl=n-1$
        \item $z=(\hat p-p_0)/\sqrt{p_0(1-p_0)/n}$
        \item Con $p_0$
        \item $z=-2{,}00$
        \item $0{,}0228$
        \item Se rechaza $H_0$
        \item $t=2{,}1429$
      \end{enumerate}
    \end{column}
    \begin{column}{0.50\textwidth}
      \begin{enumerate}\setlength{\itemsep}{0.16cm}
        \setcounter{enumi}{10}
        \item $gl=24$
        \item Se rechaza ($p=0{,}0425$)
        \item No se rechaza
        \item $z=0{,}8165$
        \item No se rechaza ($p=0{,}2071$)
        \item Sí: el IC contenía $0{,}60$
        \item El valor $p$
        \item No
        \item Se rechaza $H_0$
        \item No: solo que la evidencia es incompatible con $H_0$
      \end{enumerate}
    \end{column}
  \end{columns}
\end{frame}

\begin{frame}{Referencias}
  \footnotesize
  \begin{itemize}
    \item[{[1]}] Díaz Rodríguez, M. (2022). \textit{Estadística inferencial
      aplicada}. Universidad del Norte.
      \texttt{https://elibro.net/es/lc/uniminuto/titulos/221614}
    \item[{[2]}] Llinás Solano, H. (2014). \textit{Introducción a la
      estadística con aplicaciones en ciencias sociales}. Universidad del
      Norte. \texttt{https://elibro.net/es/lc/uniminuto/titulos/70062}
    \item[{[3]}] Contento Rubio, M. R. (2019). \textit{Estadística con
      aplicaciones en R}. Editorial Utadeo.
      \texttt{https://elibro.net/es/lc/uniminuto/titulos/220926}
  \end{itemize}

  \vspace{0.35cm}
  \begin{center}
    {\Large\color{uniAzul}\textbf{¿Preguntas?}}
  \end{center}
\end{frame}

\end{document}
